\uuid{HuGr}
\titre{Système linéaire paramétré}
\niveau{L2}
\module{Algèbre}
\chapitre{Déterminant, système linéaire}
\sousChapitre{Système linéaire, rang}
\theme{Système paramétré, discussion selon le paramètre, ensemble des solutions}
\auteur{}
\datecreate{2026-06-04}
\organisation{AMSCC}
\difficulte{3}

\contenu{
\texte{
  Soit $\lambda \in \mathbb{R}$. On considère le système linéaire :
  \[
    (S_\lambda)\;:\;\begin{cases}
      (\lambda-1)\,x + 6z = 4 \\
      y = 1 \\
      x + \lambda z = 2
    \end{cases}
  \]
}
\begin{enumerate}

\item
  \question{Écrire la matrice $A_\lambda$ et le vecteur $b$ associés au système $(S_\lambda)$.}
  \reponse{
    \[
      A_\lambda = \begin{pmatrix}\lambda-1&0&6\\0&1&0\\1&0&\lambda\end{pmatrix},
      \qquad
      b = \begin{pmatrix}4\\1\\2\end{pmatrix}.
    \]
  }

\item
  \question{Calculer $\det(A_\lambda)$ en fonction de $\lambda$ et déterminer ses racines.}
  \reponse{
    En développant par rapport à la deuxième ligne (qui contient deux zéros) :
    \[
      \det(A_\lambda) = 1\cdot\bigl[(\lambda-1)\lambda - 6\bigr]
                      = \lambda^2 - \lambda - 6
                      = (\lambda-3)(\lambda+2).
    \]
    Les racines sont $\lambda = 3$ et $\lambda = -2$.
  }

\item
  \question{Pour quelles valeurs de $\lambda$ le système admet-il une solution unique ?}
  \reponse{
    Le système admet une solution unique si et seulement si $\det(A_\lambda) \neq 0$,
    c'est-à-dire pour tout $\lambda \in \mathbb{R} \setminus \{-2,\,3\}$.
  }

\item
  \question{Pour $\lambda = 3$, le système a-t-il des solutions ?
    Si oui, déterminer l'ensemble des solutions.}
  \reponse{
    Pour $\lambda = 3$, le système devient :
    \[
      \begin{cases} 2x + 6z = 4 \\ y = 1 \\ x + 3z = 2 \end{cases}
    \]
    La première équation se réduit à $x + 3z = 2$, identique à la troisième.
    Le système est compatible, avec $y = 1$ et l'équation $x + 3z = 2$.
    En posant $z = t$ ($t \in \mathbb{R}$) comme paramètre libre, on obtient $x = 2-3t$.
    L'ensemble des solutions est :
    \[
      \mathcal{S} = \bigl\{(2-3t,\;1,\;t) \mid t \in \mathbb{R}\bigr\}.
    \]
  }

\item
  \question{Pour $\lambda = -2$, le système a-t-il des solutions ?
    Si oui, déterminer l'ensemble des solutions.}
  \reponse{
    Pour $\lambda = -2$, le système devient :
    \[
      \begin{cases} -3x + 6z = 4 \\ y = 1 \\ x - 2z = 2 \end{cases}
    \]
    De la troisième équation : $x = 2 + 2z$.
    En substituant dans la première : $-3(2+2z)+6z = 4 \Rightarrow -6 = 4$.
    C'est une contradiction, donc le système n'a \textbf{pas de solution}.
  }

\item
  \question{Pour $\lambda = 0$, résoudre le système.}
  \reponse{
    Pour $\lambda = 0$, le système devient :
    \[
      \begin{cases} -x + 6z = 4 \\ y = 1 \\ x = 2 \end{cases}
    \]
    La troisième équation donne directement $x = 2$.
    En substituant dans la première : $-2 + 6z = 4 \Rightarrow z = 1$.
    La solution unique est :
    \[
      (x,\,y,\,z) = (2,\;1,\;1).
    \]
  }

\end{enumerate}
}
